\begin{frame}[noframenumbering,plain]
    \setcounter{framenumber}{1}
    \maketitle
\end{frame}

\begin{frame}
    \frametitle{Положения, выносимые на защиту}
    \begin{itemize}
        \item В маломодовом оптическом волокне взаимное влияние одномодового и межмодового четырёхволнового смешения обусловлено перекрытием спектров усиления их общих антистоксовых компонент и устраняется фильтрацией высших мод на входе волокна.
        \item В процессе оптического полинга германосиликатного волокна излучением ближнего ИК-диапазона ширина кривых фазового синхронизма генерации второй гармоники монотонно уменьшается, а их пик сдвигается к расстройке, определяемой фазовой самомодуляцией и кросс-модуляцией. Эволюция кривых завершается при достижении предельных значений, определяемых условиями полинга, и происходит позже выхода эффективности ГВГ на постоянный уровень.
    \end{itemize}
\end{frame}
\note{
    Проговариваются вслух положения, выносимые на защиту
}


\begin{frame}
    \frametitle{Содержание}
    \tableofcontents
\end{frame}
\note{
    Работа состоит из четырёх глав.

    \medskip
    В первой главе \dots

    Во второй главе \dots

    Третья глава посвящена \dots

    В четвёртой главе \dots
}
